% Local DRF: a race on x is confined to x and to its stretch of the
% run; accesses to y, before and after, keep sequential meaning.
% Used in M10-L03.
\documentclass[tikz,border=4pt]{standalone}
\usepackage{tikz}
\input{_m10-styles.texinput}
\begin{document}
\begin{tikzpicture}[m10, x=1mm, y=-1mm, font=\large]
  % the race window (shaded band across both timelines)
  \fill[freedfill!70] (60,-4) rectangle (92,50);
  \draw[freededge, dashed, line width=0.7pt] (60,-4) rectangle (92,50);
  \node[badnote, font=\large\bfseries, anchor=south] at (76,-5) {race on \texttt{x}};

  % two domain timelines
  \node[stackcell, minimum width=26mm, minimum height=12mm, font=\large] (d1) at (4,8)  {Domain A};
  \node[stackcell, minimum width=26mm, minimum height=12mm, font=\large] (d2) at (4,40) {Domain B};
  \draw[flow] (18,8)  -- (132,8);
  \draw[flow] (18,40) -- (132,40);

  % the racy location
  \node[freedcell, minimum width=14mm, minimum height=10mm, font=\Large] (x) at (76,24) {\texttt{x}};
  \draw[badflow, line width=1pt] (68,8)  -- (x.north west);
  \draw[badflow, line width=1pt] (84,8)  -- (x.north east);
  \draw[badflow, line width=1pt] (68,40) -- (x.south west);
  \draw[badflow, line width=1pt] (84,40) -- (x.south east);

  % race-free accesses to y, before and after the race window
  \node[heapcell, minimum width=20mm, minimum height=9mm] (w) at (36,8)  {\texttt{y := 1}};
  \node[heapcell, minimum width=20mm, minimum height=9mm] (r) at (116,8) {\texttt{!y}};
  \draw[goodflow, dashed] (w.north) to[bend left=22]
    node[goodnote, above=1.5mm, font=\large, pos=0.12]{reads \texttt{1}, guaranteed} (r.north);

  % the two bounds
  \node[goodnote, anchor=north west, font=\large, align=left] at (-9,56)
    {\textbf{space}: the race cannot touch \texttt{y}\qquad
     \textbf{time}: outside the window, sequential meaning holds};
\end{tikzpicture}
\end{document}
